\subsection*{O232\quad 二次谐波的相位失配、sinc效率与准相位匹配}
无泵浦耗尽、慢变包络与平面波近似下，二次谐波振幅$A_2(z)$满足$dA_2/dz=i\kappa A_1^2 e^{i\Delta k z}$，$A_1$为常量实泵浦包络、$\kappa$实数，$\Delta k=2k_1-k_2$。
(1) 积分求长度$L$晶体的二次谐波复幅。
(2) 推导相位失配对生成强度的sinc平方规律和相干长度。
(3) 在非线性系数每个相干长度变号的准相位匹配中求首谐波有效耦合。
(4) 讨论泵浦耗尽、Manley–Rowe光子数关系与理想近似的界限。
\par\textbf{逐问解析}
(1) 若入射二次谐波为零，$A_2(0)=0$，对$z$积分得$A_2(L)=i\kappa A_1^2(e^{i\Delta kL}-1)/(i\Delta k)=i\kappa A_1^2 L e^{i\Delta k L/2}\operatorname{sinc}(\Delta k L/2)$，其中$\operatorname{sinc}x=\sin x/x$。在$\Delta k=0$处须取连续极限$A_2=i\kappa A_1^2 L$，而不是分母为零后宣布发散。
(2) 在相同模式规范下生成强度$I_2(L)\propto|A_2|^2=\kappa^2|A_1|^4L^2\operatorname{sinc}^2(\Delta kL/2)$。相位差积累到$\pi$的距离$L_c=\pi/|\Delta k|$称相干长度，超出后不同晶体段的新增二次谐波会部分相消。对精确相位匹配$\Delta k=0$，强度按$L^2$增强。
(3) 令非线性系数按符号函数$s(z)=\operatorname{sgn}[\cos(\Delta k z)]$周期翻转，其傅里叶展开在与失配相位匹配的基频处的复系数模为$2/\pi$。在跨越多个周期且忽略端部修正时，积分$s(z)e^{i\Delta kz}$的共振平均约$(2/\pi)L$，故有效系数$\kappa_{\rm eff}=(2/\pi)\kappa$，强度恢复近似$L^2$增长但相对完美天然相位匹配效率减少因子$4/\pi^2$。
(4) 如果二次谐波已经明显抽取泵浦能量，$A_1$不再常量，需联立$\chi^{(2)}$三波耦合方程；一个二次谐波光子能量$2\hbar\omega$，所以理想频率转换对应消耗两个基频光子，Manley–Rowe不变量需按光子流而非只比较电场幅值。有限群速度失配、衰减及横向衍射还会使长晶体的$L^2$增长失效。
\par\medskip
\subsection*{O233\quad 全反射界面的Goos–Hänchen位移与反射相位色散}
TE偏振平面光由非磁介质$n_1$射向较低折射率介质$n_2$，$n_1>n_2>0$，入射角$\theta>\arcsin(n_2/n_1)$，真空波数$k_0$。研究有限宽度准单色光束的反射光斑横向位移。
(1) 写出全反射区的切向波数、法向波数和倏逝常数。
(2) 求TE Fresnel反射复系数及其相位。
(3) 用驻相法推导Goos–Hänchen横向位移。
(4) 讨论临界角与掠射极限的形式发散及有限束宽修正。
\par\textbf{逐问解析}
(1) 界面切向波数守恒，$k_x=n_1k_0\sin\theta$。入射侧法向分量$k_{z1}=\sqrt{n_1^2k_0^2-k_x^2}=n_1k_0\cos\theta>0$，透射侧在全反射区为$k_{z2}=i\kappa$，$\kappa=\sqrt{k_x^2-n_2^2k_0^2}>0$；因此低折射率侧场沿法线指数衰减而无传播法向功率。
(2) 对TE偏振，Fresnel反射系数$r_s=(k_{z1}-k_{z2})/(k_{z1}+k_{z2})=(k_{z1}-i\kappa)/(k_{z1}+i\kappa)$，故$|r_s|=1$，相位可取$\phi_r(k_x)=-2\arctan(\kappa/k_{z1})$。虽反射率为1，但相位随入射波数变化，形成反射波包空间移动的机制。
(3) 将入射束作围绕中心$k_x$的角谱积分，反射振幅额外含$e^{i\phi_r(k_x)}$。对窄角谱作驻相，光斑位移$D=-d\phi_r/dk_x$。利用$dk_{z1}/dk_x=-k_x/k_{z1}$、$d\kappa/dk_x=k_x/\kappa$，得到$d[\arctan(\kappa/k_{z1})]/dk_x=k_x/(\kappa k_{z1})$，故$D=2k_x/(\kappa k_{z1})$，单位是长度且为正。
(4) 当$\theta\downarrow\theta_c$，$\kappa\to0$使平面波公式的$D$形式发散；当$\theta\to\pi/2$，$k_{z1}\to0$也会产生形式发散。真实有限宽度束在临界角附近有宽角谱，不能将反射相位仅作一阶Taylor展开，需用完整角谱积分；束宽、偏振和材料吸收都会使位移有限或改变峰位。
\par\medskip
